% Part: sets-functions-relations
% Chapter: sets
% Section: basics

\documentclass[../../../include/open-logic-section]{subfiles}

\begin{document}

\olfileid{sfr}{set}{bas}
\olsection{د غړو له مخې برابري}

\emph{سټ} د څيزونو يوه ټولګه ده چې د يوه څيز په توګه ورته کتل کېږي۔ هغه څيزونه چې سټ ترې جوړ وي، د سټ \emph{عناصر} يا \emph{غړي} بللے شي۔ که $x$ د يو سټ~$A$ !!a{element} وي، $x \in A$ ليکو؛ که نۀ وي، $x \notin A$ ليکو۔ هغه سټ چې هېڅ !!{element}s نۀ لري، \emph{تش} سټ بللے شي او په ``$\emptyset$'' ښودلے کېږي۔

\begin{explain}
دا مهمه نۀ ده چې سټ څنګه \emph{ټاکو}، د هغۀ !!{element}s څنګه \emph{ترتيبوو}، يا په رښتيا د هغۀ !!{element}s \emph{څو ځله} شمارو۔ يوازې دا مهم دي چې د هغۀ !!{element}s کوم دي۔ دا خبره په لاندې اصل کښې راټولوو۔
\end{explain}

\begin{defn}[د غړو له مخې برابري]
  که $A$ او $B$ سټونه وي، نو $A = B$ هله او يوازې هله وي چې د~$A$ هر !!{element} د~$B$ هم !!a{element} وي، او برعکس۔
\end{defn}

د غړو له مخې برابري د ځينو نښو د کارولو جواز راکوي۔ په عمومي توګه، چې کله ځينې څيزونه $a_{1}$، \dots، $a_{n}$ ولرو، نو $\{a_{1}, \dots, a_{n}\}$ \emph{هم هغه} سټ دے چې !!{element}s ئې $a_1, \ldots, a_n$ دي۔ په ``\emph{هم هغه}'' ټينګار کوو، ځکه چې د غړو له مخې برابري راته وائي چې داسې سټ يوازې \emph{يو} کېدے شي۔ په حقيقت کښې، د غړو له مخې برابري د لاندې خبرې جواز هم راکوي:
  \[
    \{a, a, b\} = \{a, b\} = \{b,a\}.
  \]
دا هم هغه خبره څرګندوي چې د سټونو په پام کښې نيولو پر وخت، د هغو د \psOblique{!!{element}s} ترتيب يا دا چې څو ځله ښودل شوي دي، مهم نۀ دي۔

\begin{tagblock}{novice}
\begin{ex}
چې کله د څيزونو يوه ډله ولرئ، هغه په يو سټ کښې راټولولے شئ۔ د بېلګې په توګه، د رېچرډ د وروڼو او خوېندو سټ يو داسې سټ دے چې يو کس لري، او هغه د $S=\{\textrm{Ruth}\}$ په بڼه ليکلے شو۔ له $4$ څخه د کوچنيو مثبتو صحيح عددونو سټ $\{1, 2, 3\}$ دے، خو د $\{3, 2, 1\}$ يا آن د $\{1, 2, 1, 2, 3\}$ په بڼه هم ليکل کېدے شي۔ د غړو له مخې د برابرۍ له مخې، دا ټول هم يو سټ دے۔ ځکه چې د $\{1, 2, 3\}$ هر !!{element} د $\{3, 2, 1\}$ (او د $\{1, 2, 1, 2, 3\}$) هم !!a{element} دے، او برعکس۔
\end{ex}
\end{tagblock}

ډېر ځله به يو سټ د يو داسې خاصيت په وسيله ټاکو چې د هغۀ !!{element}s ئې په ګډه لري۔ د دې دپاره به دا لنډه نښه کاروو: $\Setabs{x}{\phi(x)}$، چې پکښې $\phi(x)$ د هغه خاصيت استازيتوب کوي چې~$x$ ئې بايد ولري، څو د سټ په \psOblique{!!{element}s} کښې وشمېرل شي۔

\begin{tagblock}{novice}
\begin{ex}
زمونږ په بېلګه کښې، $S$ په دې بڼه هم ټاکلے شو:
\[
S = \Setabs{x}{x \text{ د رېچرډ ورور يا خور دے}}.
\]
\end{ex}
\end{tagblock}

\begin{tagblock}{math}
\begin{ex}
يو عدد هله او يوازې هله \emph{کامل} بللے شي چې د خپلو حقيقي قاسمونو له مجموعې سره برابر وي (يعنې هغه عددونه چې دا عدد پرې پوره وېشل کېږي، خو پخپله له دې عدد سره برابر نۀ وي)۔ د بېلګې په توګه، $6$ کامل دے ځکه چې حقيقي قاسمونه ئې $1$، $2$ او~$3$ دي، او $6 = 1 + 2 + 3$۔ په حقيقت کښې، $6$ له $10$ څخه يوازينے کوچنے مثبت صحيح عدد دے چې کامل دے۔ نو د غړو له مخې د برابرۍ په کارولو سره وئيلے شو:
\[
	\{6\} = \Setabs{x}{x\text{ کامل دے او }0 \leq x \leq 10}
\]
ښي لاس ته نښه داسې لولو: ``د هغو $x$ ګانو سټ چې $x$ کامل وي او $0 \leq x \leq 10$ وي''۔ دلته برابري دا خبره تائيدوي چې د سټونو په پام کښې نيولو پر وخت، د هغو د ټاکلو طريقه مهمه نۀ ده۔ او په لا عمومي توګه، د غړو له مخې برابري ډاډ راکوي چې تل د هغو $x$ ګانو يوازې يو سټ وي چې $\phi(x)$ وي۔ نو د غړو له مخې برابري دا جواز راکوي چې $\Setabs{x}{\phi(x)}$ ته د هغو $x$ ګانو \emph{هم هغه} سټ ووايو چې~$\phi(x)$ وي۔
\end{ex}
\end{tagblock}

د غړو له مخې برابري راته د دې ښودلو يوه طريقه راکوي چې سټونه برابر دي: د $A = B$ د ښودلو دپاره، وښايئ چې هر کله $x \in A$ وي نو $x \in B$ هم وي، او هر کله $y \in B$ وي نو $y \in A$ هم وي۔

\begin{prob}
ثابت کړئ چې له يو څخه زيات تش سټونه نۀ شي کېدے، يعنې وښايئ چې که $A$ او $B$ داسې سټونه وي چې هېڅ !!{element}s نۀ لري، نو $A = B$۔
\end{prob}

\end{document}
